\section{Keeping company choices and their consequences together}
\label{sec:company-family-language}

The fixed collection amounts and the 90-percent rule answer different questions
about the same company opening. An analyst comparing them needs each choice's
cash, receivables, profit and later balances together. Selecting favorable
amounts independently would describe a company that no stated choice produces.
A named scenario family therefore retains the entire input path and result of
each member, including its unresolved years.

The bounded family contains the baseline, two price revisions, the unknown
payroll path, the supplier-payment failure and the policy-generated path. Four
members complete their accounting; two remain partial. Every member has the
same historical source, opening record, entity, book, currency and reporting
periods. Opening claims, asset memory, drawn facility balance and debt principal
must also agree. Future instructions may differ. These checks establish the
common declared opening; they do not reconstruct an unobserved real history.

The text places an existing company path inside a named member block. Typed
constructors embed the same path record. A family query returns each member's
complete report vector for the requested year and retains its path and result
identities. If a member never reached that year, its response says why the
report is unavailable. The full family result also retains the preceding and
following years, so a later amount cannot be detached from the path that
generated it.

These are authored exogenous scenarios. Six names imply neither equal
probabilities nor exhaustive coverage of possible futures. Two declarations
that produce identical results can remain separately named assumptions without
acquiring extra probability mass. An empty authored selection means that the
author supplied no members; it does not establish that no history is feasible.
Historical reconstruction and justified probability laws retain their separate
requirements from Chapter~\ref{ch:math}.

\paragraph{A policy expression has units and dependencies.}
The collection fraction can now be written as a calculation rather than supplied
only as a final number. In the ordinary scenario it is one minus a declared
discount of one tenth. A second branch supplies four fifths. Both branches use
two equal instalments. The guard is an explicit Boolean assumption: it does
not infer whether customers will behave normally.

The expression graph contains exact rational or Boolean inputs, literals,
references and pure arithmetic, comparison and conditional operations. Text
and typed construction preserve the same graph. Static analysis checks all
references, operator arities, cycles and types, including both conditional
arms. The resulting fraction and instalment shares pass through the existing
settlement policy's domain and claim-allocation checks. The expressions do
not create new collection dates or decide which contracts should exist.

Numeric types carry units and an explicit context. A quantity of goods times
a price in dollars per good gives dollars; a quantity cannot be added to cash.
Addition and comparison require matching dimensions and context. A
context-free dimensionless fraction can scale a company amount. Multiplication
does not silently convert currencies or combine different entities. A ratio
of contextual amounts retains that context and cannot automatically become
an unrestricted policy parameter.

For Northstar's two revised invoices, the total amount is 384. The ordinary
fraction of nine tenths collects 345.6; four fifths collects 307.2. The latter
therefore reduces cash by 38.4 and increases receivables by the same amount,
with unchanged profit. The exact arithmetic agrees with the generated path's
full accounting result. The example establishes the consequence of the two
stated collection assumptions, without estimating which assumption is more
likely.

A known guard evaluates only its selected arm, after both arms have passed
static checks. An unknown guard leaves the policy output unresolved; it is
not false. This bounded evaluator does not prove that equal-looking arms are
symbolically equivalent. A valid but undefined calculation in an unselected
arm is not evaluated, while division by known zero in a selected calculation
is invalid. Missing declarations differ from declared unknown values. The
former are invalid references; the latter remain visible dependencies.

Each output records all declared input dependencies and the subset actually
evaluated. Thus an unselected branch's assumptions remain inspectable even
when they do not determine the current value. Logical \code{and} and
\code{or} conservatively propagate unknown operands in this subset; they do
not introduce implicit short-circuit or SQL missing-value semantics. General
functions, date arithmetic, source maps, richer types and endogenous event
scheduling remain language work.

\paragraph{Reuse a calculation without changing its meaning.}
The collection example can express the ordinary fraction through a reusable
function: given a dimensionless discount \(d\), return \(1-d\). A second
function selects the ordinary or stress fraction from the explicit Boolean
assumption; a third returns the equal instalment share \(1/2\). Calling
these functions with the existing inputs yields the same \(9/10\), \(4/5\)
and \(1/2\) values as the earlier expression graph. Their reuse changes
how an assumption is authored, not how it is estimated.

The optional \code{CompanyFunctions/1} declaration names each function's
ordered, explicitly typed parameters, return type, local expression nodes
and output. A call supplies argument-node identities in parameter order.
Function bodies can read their parameters and literals; they cannot capture
an undeclared company input. Expansion gives each call fresh local identities
and lowers its operations to the existing expression evaluator. Parameter
and return types must agree exactly, including units and company context.
Every definition receives static checks, even if no selected output uses it.

Evaluation follows the existing graph's demand for values. An unused argument
does not have to be evaluated, although its references and types must be valid.
For example, a function that ignores one argument can return a known constant
even when that argument contains a well-typed division by zero. If the result
uses that division, it is invalid. Likewise, the stress collection branch can
remain known when its unused discount is unknown. Each call retains all
declared argument dependencies separately from the dependencies that reach its
result; output evaluation records the inputs actually used.

The checked comparison preserves complete three-year results for the ordinary
and stress policies. Moving from \(9/10\) to \(4/5\) reduces cash by
\(192/5\) and raises receivables by \(192/5\), with unchanged profit.
An unknown selected discount or guard prevents policy materialization. A
negative generated fraction is rejected by the existing collection policy's
domain checks. Compilation supplies neither a collection assumption nor a
missing fact.

The bounded reference accepts at most 32 functions, 16 parameters per
function and 256 source nodes per scope. Call depth is at most 16; both
local and expanded dependency depth are at most 64. Validation and final
expansion together permit at most 2,048 emitted nodes and 2,048 scope visits,
including calls that merely pass an argument through. Serialized requests
are limited to 262,144 UTF-8 bytes. These are explicit convenience limits,
not measured scalability claims. Cycles, recursion and exceeded limits are
refused; there is no implicit numerical solution or event scheduling.

The function request, compiler and generated legacy graph retain separate
identities and definition/call origins. Their information evidence does not
inherit the earlier policy's dates. This increment provides typed declarations;
matched textual function syntax, parser spans, richer types and module dispatch
remain further language work. Its 23 focused methods, 79 predecessor checks
and 46 policy/accounting comparisons are retained in
\path{data/accounting-master/company-functions-20261025/}.


\paragraph{Locate each declaration in the source that was supplied.}
The matched text extension now expresses these same bounded functions. For
example, the ordinary collection fraction is written as follows:
\begin{lstlisting}[basicstyle=\ttfamily\scriptsize,breaklines=true]
function complement;
parameter discount '{"context":null,"dimensions":{},"kind":"number"}';
returns '{"context":null,"dimensions":{},"kind":"number"}';
param_ref d discount;
literal one '{"context":null,"dimensions":{},"kind":"number"}' 1/1;
node r sub one d;
result r;
end_function;
\end{lstlisting}
The quoted JSON type denotes a number with no dimensions and no company
context. The parameter reference binds the local name \code{d} to the declared
discount; \code{r} subtracts it from one. The result declaration identifies
the returned node, and \code{end\_function} closes the function. In the
surrounding program, \code{call regular complement discount;} supplies the
main program's discount node and names the resulting call \code{regular}.
Input declarations and the other functions remain necessary for the full
collection policy. The displayed wrapping does not introduce extra source
lines: each declaration occupies one physical line and ends in a semicolon.

Function parameters precede the return type and body; a result ends the body.
Blocks cannot nest, but a call can name a function declared later. Main inputs,
nodes and outputs use the earlier expression notation, within a
\code{CompanyFunctions/1} program ending in \code{end\_program}.
The parser rejects duplicate names and JSON keys, incomplete blocks, unsupported
declarations and semantic errors reported by the pinned typed compiler.
Quoted command-looking characters are text: parsing performs no shell
expansion, arithmetic evaluation or event generation. A validly typed division
by zero can therefore compile and later fail when its value is requested.

Each accepted declaration records its path within the typed request, physical
line and exact original UTF-8 byte interval. Byte offsets are zero-based and
end-exclusive; line and Unicode-codepoint columns are one-based, with the ending
column exclusive. The slice includes the terminal semicolon and excludes
surrounding spaces, tabs and the line ending. Its hash and the complete source
hash bind these locations to the supplied bytes. This distinction matters when
non-ASCII names use several bytes per character, or a source uses CRLF rather
than LF. Both line-ending forms are accepted without Unicode normalization.

Canonical rendering uses LF and deterministic quoting. It reproduces the same
typed request but can have a different source hash and different locations.
Consequently, original locations are never assigned to re-rendered text.
The generated expression's function and node origins can be joined to these
declarations; token-level debugging and runtime exception locations remain
separate work. Source locations also supply no evidence of when a policy was
written, published or used.

The text frontend bounds source size to 524,288 bytes, each physical line to
16,384 bytes and accepted declarations to 12,000 before the typed compiler's
own limits apply. These are scoped engineering limits. Eighteen focused test
methods, 23 typed-function regression methods, 79 expression checks and 285
source/financial comparisons cover six policy cases, preserving complete
three-year results where available and unresolved or rejected paths elsewhere.
The exact examples, declaration slices and financial comparisons are retained
in \path{data/accounting-master/company-function-text-20261025/}. The checked
extension still does not supply dates, polymorphism, module dispatch or
endogenous economic decisions.


\paragraph{Let the next sale depend on the stock that remains.}
A reusable calculation becomes a company decision when it reads the current
book and proposes an economic event. Consider a hypothetical company holding
two units of a homogeneous, divisible inventory lot. The lot has a carrying
amount and tax base of USD~20, and cash is USD~452. Management's authored rule
keeps half a unit in reserve and sells the lesser of declared demand and stock
above that reserve. Demand, the reserve and the price are scenario assumptions;
the example constructor requires them to be nonnegative. The rule applies only
when its separate Boolean switch is enabled.

With demand of one unit on each of three successive decision dates and a price
of USD~15 per unit, the first decision sells one unit. The second reads the
remaining unit and sells only half. The third reads the remaining half unit
and proposes no sale. Total proceeds are USD~22.50, released inventory cost is
USD~15 and profit before any subsequent tax assessment is USD~7.50. Cash becomes
USD~474.50, while the retained half unit carries USD~5. Reusing the opening stock
for all three decisions would give the wrong quantities. These exact fractions
provide an independent check on the state read at each decision.

The bounded reference \texttt{StateDecisionEvent/1} separates this choice from
accounting for its consequences. It receives the complete conditional annual
book, its content hash, a decision date, a typed function program and explicit
input bindings. A binding can read the selected lot's current quantity or the
book's cash balance, or supply a scenario value with provenance and an
availability date. Quantity, currency and company context must agree with the
program's declared types. Every program input starts as an unknown placeholder;
the binding determines which value the rule may read.

A value dated after the decision, or with unknown availability, stays unknown
for evaluation while its original record remains inspectable. The model itself
also needs a declared availability date. These dates express the authored
scenario; they do not establish when a real company knew a fact. The caller
must further state that no intervening economic event is omitted between the
snapshot and the decision. The reference neither verifies that premise against
disclosures nor discovers event dates. Decision and sale share one calendar
date within the current fiscal period; explicit order distinguishes decisions
on the same day.

The program supplies three outputs: whether to act, the units sold and the
proceeds. The compiler checks the entire program before evaluation. A false
decision guard leaves the book exactly unchanged and need not read an unused
future price. An unknown guard leaves the decision unresolved. Only a true
guard requires the units and proceeds; unknown values remain unresolved, while
invalid arithmetic refuses the proposal. This distinction prevents missing
information from silently becoming a zero sale.

A known proposal still needs the existing sale method's six findings:
proportional allocation selected, homogeneous units, divisible units, unchanged
measurement basis, completed performance, and tax base equal to carrying amount.
The decision rule supplies none of these findings by itself. Application
recomputes the complete proposal against the same state and request, then
dispatches it to the unchanged native sale transition. A stale snapshot or an
altered proposal refuses even if the caller updates its hash. Unknown findings
remain unresolved, false findings remain unsupported, and a refused proposal
is retained separately from its accounting outcome.

In the checked single and repeated cases, the proposed operations, complete
native transitions and financial statement fields equal those produced by
manually authored sales. The report describes the recorded decision prefix and
retains the annual interface's unresolved questions. A sequence stops at its
first refusal, preserving earlier successful decisions, no-actions and remaining
requests. Its decision cursor can advance through a no-action even though the
book's last event date does not change; a later request cannot move behind that
cursor. The reference accepts at most 64 decisions and 32 input bindings per
decision. These are engineering limits, not a capacity claim.

Twenty-eight focused test methods, 14 native-accounting regression methods,
23 function regression methods and 64 exact comparisons across eight cases are
retained in \path{data/accounting-master/state-decision-events-20261025/}.
They check this composition of an authored rule and a conditional accounting
method. The example leaves demand estimation, general financing decisions,
simultaneous choices, event scheduling and actual historical execution open.
No database commit, complete fiscal event population or forecast accuracy
follows from the successful sequence.


\paragraph{Bind a scheduled decision when its date arrives.}
The preceding sequence fixes the state read by each decision, but its caller
must supply the corresponding future book hashes. An authored simulation should
instead supply its starting book and a calendar, then calculate those later
states as it proceeds. The bounded \texttt{StateDecisionSchedule/1} extension
does this by materializing each decision request only at its turn. Its expected
state hash is the hash of the current book, after earlier successful events.

A schedule supplies a named rule, explicit priority and either a list of
calendar dates or a start date, exclusive end date and positive integer number
of days between occurrences. For example, daily decisions from 2 January 2027
up to, but excluding, 5 January produce the same three decisions as the explicit
dates 2, 3 and 4 January. Calendar arithmetic includes leap days. It introduces
no month-end adjustment, business calendar, timezone conversion or contractual
day-count convention. All occurrences must fall within the current fiscal
period and no earlier than the book's last event.

Different decisions on the same day need distinct explicit priorities. An
unresolved tie stops scheduling before economic execution; rearranging the input
array does not choose an order. Each decision and sale receives a stable identity
derived from the declared root and schedule names and its occurrence number.
The complete schedule also retains its own content identity, so an equivalent
list of dates and recurrence can produce identical execution while remaining
different authored documents. Every declared policy is statically checked
before the first event, including a later policy that an earlier refusal might
otherwise prevent the runner from reaching.

The checked calendar reproduces the manually assembled stock/reserve sequence:
one unit sold, then half a unit, then no action. Cash is USD~474.50. The selected
lot retains USD~5 of carrying amount, while the inventory control totals USD~65:
the opening USD~80 also included USD~60 outside that lot. With a zero reserve,
the same rule sells both units over two decisions and cash reaches USD~482;
the next decision observes depletion and takes no action. Distinguishing a
selected lot from the company's inventory control is necessary when checking
the full balance sheet.

An unresolved or refused decision stops the sequence with the exact current
book, earlier outcomes and pending calendar retained. The stop record includes
the materialized request, its state hash and the proposed event where available.
A no-action advances the decision cursor without changing the book's event
date. The scheduler accepts up to 32 schedules and 64 total decisions; the
existing request checker also bounds serialized size, nesting and structure.
These limits keep the reference finite and do not establish production capacity.

Eighteen focused scheduler methods, 28 direct decision regression methods and
49 exact comparisons across five cases preserve manually assembled requests,
complete native outcomes and all financial report fields. The evidence is in
\path{data/accounting-master/state-decision-schedule-20261025/}. The calendar
retains the explicit premise that no external economic event is omitted.
Combining decisions with separately supplied financing, tax, purchase or period
events requires those events' own timing and accounting rules. General calendar
generation, historical completeness and durable restart remain separate work.


\paragraph{Declare the company book before applying its methods.}
The preceding native examples use a fixed company context and chart of accounts.
A reusable calendar does not remove that restriction. A company may keep its cash
under account \texttt{101}, another under \texttt{bank.alpha}; their contract
records may also have different structures. The optional
\texttt{CompanyBook/1} reference makes those choices explicit before a method
reads balances or proposes an entry. It declares the entity, book, currency,
conditional case, period, account classifications and selected cash account.
Balances use the existing debit-positive convention and sum to zero. Each scalar
in the method's memory has a declared type and an address within a closed schema.

The payment example separates the accounting relation from those local names.
One synthetic company maps the cash and liability roles to
\texttt{bank.alpha} and \texttt{note.alpha}, with the outstanding amount under
\texttt{loan}. Another maps the roles to \texttt{101} and \texttt{230}, with the
amount under \texttt{contracts/obligation}. Both start with cash of USD~100 and an
outstanding liability of USD~30. A payment of USD~3.50 credits cash and debits
the liability by the same amount, leaving cash of USD~96.50 and a liability of
USD~26.50. The contract record falls to USD~26.50 and its cumulative payments
rise to USD~3.50. A second payment of USD~26.50 extinguishes the liability and
leaves cash of USD~70. Other accounts and unrelated contract records are unchanged.
The example assumes sufficient cash and a positive payment no larger than the
outstanding amount; the generic book does not impose a universal funding policy.

\texttt{CompanyBookMethod/1} declares an injective mapping from account roles to
the book's accounts, prerequisite predicates, monetary calculations, memory
replacements and postconditions. Bindings select declared account balances,
typed memory values or supplied event inputs. Their company, dimensions and
scale must agree exactly with the method's input declarations. A dollar amount
cannot become an inventory quantity merely because both are rational numbers.
The existing monetary interpreter derives the journal and separate gross cash
flows. The memory interpreter then replaces the specified scalar fields; this
profile rejects dynamic paths, new collection members and container writes.
Nested prerequisite and monetary definitions carry the enclosing method's
identity and version. A literal memory replacement is interpreted in the
destination's declared type.

\texttt{CompanyBookTransition/1} checks the entire request before assessing its
prerequisites. It derives monetary effects, stages memory changes and constructs
a new book from the same starting state. Only then does it evaluate the
postconditions against that proposed book. In the payment case, a postcondition
requires the credit balance of the liability account to equal the outstanding
contract amount. A failed postcondition returns no applied book, although the
staged calculation remains available for diagnosis. An unknown applicability
finding remains unresolved; a false finding can make the chosen method
unsupported. An inactive guarded prerequisite does not, by itself, cancel the
economic event: a separate decision determines whether an event should occur.

The request identifies both the starting book and the complete method by digest.
Application recomputes the proposal and compares its canonical JSON representation,
including Boolean and numeric types. Changing a saved journal, memory value or
authorization flag therefore cannot pass merely by replacing its digest. Event
identities cannot repeat, dates cannot precede the current book, and each event
must fall within the declared period. This is an atomic construction of a
provisional value; the separate posting-authorization and durable-storage rules
still govern any admitted or persisted transaction.

Five comparisons connect this reference to the retained native methods. Two
current-tax payments and three complete-lot sales reproduce their native journals,
gross flows and entire resulting states through explicit mappings. The mappings
retain every native field outside the declared changes; both payment cases also
reproduce the complete native financial reports. These checks concern the selected
synthetic cases. General reporting and closing, source admission, dynamic contract
collections, currency conversion and the broader corporate package still require
their own methods. The declared book supplies a common place to compose those
methods without weakening the older profiles' invariants.

\paragraph{Bind a named method to the event that actually occurred.}
The declared book separates company-specific accounts from a method's financial
roles. Selecting that method also requires an unambiguous connection to the event.
\texttt{CompanyMethodRegistry/1} names the available data-only method definitions,
their schemas, expected account classifications and input types. Each entry
identifies its complete method body. The optional \texttt{CompanyMethodCall/1}
selects one entry under the book's company context and checks every entry's
structure before evaluating the selected method's prerequisites. An invalid unused
entry therefore cannot remain hidden until a later event happens to select it.

Some inputs are supplied facts; others describe the event itself. Keeping these
sources distinct matters for a sale after a tax close. Suppose the event is dated
15 June, but a separately supplied date says 1 September. A predicate comparing
that supplied date with a 1 July close could accept the wrong timing. The named
binder instead derives designated date and identity inputs from the event's own
metadata. It rejects a supplied input with the same name. In the checked case,
the 15 June sale consequently fails the closed-period prerequisite. A genuine
change to the event date changes the predicate's input with it.

The binder materializes the ordinary company-book transition request, including
the current state and selected method identities. The existing prerequisite,
monetary and memory engines then determine the result. Two company payment layouts
and a complete-lot sale reproduce their explicitly authored requests and entire
provisional transitions; unknown and false applicability findings retain their
different consequences. Repeated calls read the resulting book, so the next
payment uses the reduced outstanding amount. Application recomputes the complete
named call and uses the same type-sensitive JSON comparison as the underlying
transition.

The current registry supports a directly dated event. Assessment intervals,
measurement dates and period carry are different clock families and remain
unsupported here. The registry also leaves native applicability to an explicit
adapter. For example, a reusable inventory-sale adapter must validate all lot
allocations, disposition and tax-base conditions, the selected sale method and the
native close date, as well as preserve every unmapped field. Matching one native
journal cannot establish those conditions. The selected comparisons establish the
named binding mechanism; they do not turn the retained test projections into
general native-method admission.

\paragraph{Carry the entire native book through the selected sale.}
A complete state check precedes this sale. The selected
\texttt{NativeCompanySale/1} adapter extends the earlier test projection to the complete retained native whole-lot sale domain. It validates
the native state's structure and accounting controls before importing every
existing lot into a declared company book. The native account classifications
and company context supply an explicit mapping. Fields outside the sale remain
bound to the original state and are restored when the proposed book is exported.
The adapter then validates the entire restored native state again.

Preserving every lot makes the comparison discriminating. In a two-lot example,
the inventory account holds USD~20: USD~15 belongs to the selected lot and USD~5
to a second lot. Selling the selected lot for USD~40 raises cash from USD~454
to USD~494, reduces inventory to USD~5 and raises cost of sales from USD~56 to
USD~71. The second lot is unchanged. Selling it later for USD~8 leaves cash of
USD~502 and no inventory. The declared method reconciles the inventory account
with the carrying amounts of all existing lots and other inventory before and
after each sale; it replaces only the selected lot's units, carrying amount,
tax base and disposition fields.

The generic monetary and memory interpreters compute these changes through the
named method binder. The native transition remains the comparison calculation:
a test disables it while the adapter still proposes and applies the sale. Across
23 accepted and refused cases, complete native states, journals, gross cash flows
and refusal reasons agree. Unknown performance or proceeds remain unresolved;
false performance and a fractional sale remain unsupported. Native cash-funding
failure, duplicate identities, backdating, tax-close restrictions and inconsistent
lot or control balances retain their native consequences.

This adapter preserves the fixed native profile's domain rather than inferring
accounting applicability from the general book's type checks. It creates no new
lot and supports only a complete-lot sale. Source admission, partial disposals,
new collection members and a general company's full reporting package remain
separate work. Its contribution is a checked route through the common book that
preserves the native conditions and every field outside the selected transaction.

\paragraph{A payment can change the decision that follows it.}
Separately supplied economic events can alter the state read by a scheduled
rule. The bounded \texttt{MixedEventTimeline/1} joins directly dated native
payments and inventory sales with the preceding decision calendar. It preserves
the native accounting operation and materializes each decision against the book
at its turn. Every action has an explicit date and priority; a same-day tie
remains unresolved rather than inheriting the order of the input array.

An order-sensitive example starts with cash of USD~452 and tax payable of
USD~10.25, split between carried vintages of USD~5.25 and USD~5. Management's
authored rule permits a one-unit inventory sale at USD~15 only if cash is below
USD~450. Paying USD~4 from the first tax vintage before the decision reduces
cash to USD~448 and enables the sale. Cash then becomes USD~463, total inventory
falls from USD~80 to USD~70, and the selected lot's carrying amount falls from
USD~20 to USD~10. The sale contributes USD~5 of profit; the payment reduces cash
and payable without creating a new expense.

Reverse the explicit same-day order and the decision sees cash of USD~452,
so it takes no action. The later payment leaves cash of USD~448 and inventory
of USD~80. In both cases tax payable is USD~6.25, with vintages of USD~1.25 and
USD~5. These differences follow from the specified threshold and order, not an
estimated company policy. The comparison checks the complete native journals,
state histories and financial reports as well as these independent controls.

The initial adapters accept allocated payments of carried tax vintages and the
existing proportional inventory sale. Their declared schedule date must match
the date inside the native event. Structural checks cover every supplied entry
before execution, but remaining cash, stock, payable and accounting findings
are assessed against the book at the relevant action. A later unknown vintage
or unavailable authorization therefore preserves the earlier successful prefix
and the exact failed event. An unresolved decision after a successful payment
retains that payment; it is not rolled back merely because the subsequent sale
could not be evaluated.

Assessment and period-boundary operations need different treatment. A current
tax assessment names an interval, a deferred measurement binds to a particular
state, and a carry changes the fiscal period. The mixed reference leaves these
operations, and nested return/lease envelopes, unsupported rather than assigning
them an arbitrary event date or rewriting a stale measurement. This is an
explicit implementation gap in the integrated timeline; the separate native
methods retain their existing scope.

Seventeen focused methods, 18 scheduler regression methods and 61 exact
comparisons across six cases are retained in
\path{data/accounting-master/mixed-event-timeline-20261025/}. The runner bounds
the combined population to 64 actions, including at most 32 supplied native
entries, and retains the current book, successful outcomes, failed action and
remaining timeline on refusal. Its closed-path premise covers both supplied
events and decisions. General financing, richer calendars, historical
completeness and production durability remain open.


\paragraph{Available inputs do not establish a saved past calculation.}
The path's information date cannot also date the separate policy body or its
materializer. The synthetic extension gives the policy public availability on
10 March 2025 and the method on 12 March, both at calendar-day precision.
Even if the parent path is already available, the public dependency selection
cannot pass until 13 March. The query binds observations to the exact policy
and method bytes, then recomputes the generated path and its accounting result.
Changing the rule or materializer prevents reuse of the old observations.

That result is a consequence calculated from the eligible synthetic inputs.
It receives no invented authoring, publication, ingestion or acceptance date
of its own. Showing that a calculation actually existed before a cutoff needs
separate historical execution evidence. Eligible dependencies can also produce
a partial accounting result: availability does not supply missing payroll,
closing memory or a third-year opening.

The expression program and its assumptions introduce further dependencies.
Their information evidence is not inherited from the directly supplied
90-percent policy. The current expression and family examples therefore retain
archival scenario meaning, with historical, temporal and forecast admission
unestablished. Their draft content hashes use the existing company engine's
encoding, including its non-ASCII escaping. Import into the separately defined
portable record encoding requires an explicit identity mapping;
Section~\ref{sec:company-module-language} supplies that mapping for these
legacy drafts and retains the separate verification levels.

\paragraph{Interrupted work differs from incomplete accounting.}
Suppose evaluation stops after two members. A checkpoint records those exact
member results and the next member to evaluate, bound to the original family,
source and method. Resume validates the stored prefix by replay and evaluates
the requested number of new members. The ordering and identities prevent a
member from being skipped or counted twice. Different budget splits reproduce
the same final family.

The unknown-payroll and supplier-failure members are finished evaluations with
partial financial results. They are not unfinished computational work. Once
stored, they remain part of the completed prefix, and recovery continues with
the next member. This distinction is necessary to avoid repeatedly evaluating
a known unresolved case while silently excluding it from the family. Recovery
inside a member's year or event sequence remains outside this checkpoint.

The scheduling budget counts newly evaluated members. Validating a saved
prefix also takes work, recorded separately; the budget is not a claim about
total runtime. A changed source, model or method requires a new operation,
and a rehashed false prefix fails trusted replay. Neither changed semantics
nor a schema migration can be disguised as resume.

Publication of the checkpoint has its own work owner and generation. One
reference transition publishes the new checkpoint and immutable request
decision together. If acknowledgement is lost, either nothing was published
or the complete transition occurred. Retrying the exact request publishes it
in the first case and returns the saved response in the second. That response
remains the same after a later request advances the work head. Reusing the
identity with different request bytes conflicts; a new request with a stale
head does not advance it.

The reference rejects a world with only part of the checkpoint, decision and
generation update. It checks these publication meanings through explicit
before-and-after states and complete replay, without establishing physical
storage durability or changing the company's accounting state. A persistent
implementation, efficient recovery and compatible migrations remain distinct
engineering obligations.

The four extensions are retained under \path{data/accounting-master/}:
\begin{quote}\small
\path{company-policy-information-20261023}\\
\path{company-family-20261023}\\
\path{company-expressions-20261023}\\
\path{company-recovery-20261023}
\end{quote}
They connect bounded authoring, information and recovery behavior to the same
company example. Broader corporate and
sector methods, actual source admission and independent implementation
evidence remain necessary for the full financial-statement language.

\section{Preserving a model before extending its accounts}
\label{sec:company-module-language}

The scenario family now keeps each company's instructions and consequences
together. Exchanging those results with another implementation raises two
further questions. Can it identify the same input despite a different text
encoding? Can it add a new accounting method without creating another cash
balance or losing the company's tax history? The first is a representation
question. The second requires a richer account of company state.

\paragraph{Different bytes can express the same instruction.}
A company name containing an accented character may be written directly in
UTF-8 or through a JSON escape. Both decode to the same sequence of characters,
but their original files differ. Keeping the exact source file and identifying
the decoded model therefore require separate identities.

The legacy import retains the original bytes and their hash, the declared local
decoder and profile, and the existing company draft's decoded identity. A
separate portable content record binds the canonical decoded value. Equivalent
escaping or whitespace can share that content record while retaining distinct
original imports. Different character sequences are not normalized into one
name. The exchange itself changes when its original bytes change.

A field extracted from a saved result is not passed off as an original file.
Its import identifies the parent and exact JSON pointer, and validation checks
that the value is present there. Family-member records expose names, paths,
results, contexts and annual outcomes; expression records connect the program,
policy, selected parent, generated child and declared dependencies. Recovery
records retain the same family and method with the completed member prefix.
Missing dependencies and wrong reference kinds prevent import. This is an
explicit reference to legacy drafts, rather than a conversion of every field
into the full portable numeric and type language.

Identity checks alone cannot establish that cash was calculated correctly.
One adverse example changes a cash amount and recomputes the affected hashes,
including the saved recovery result. The resulting packet is internally
consistent and passes content binding. Repeating the calculation with the
trusted local engine rejects it. The two verdicts remain distinct: a complete
dependency graph preserves what was supplied, while semantic replay checks
what the selected method computes. Neither supplies an actual source's
publication date, proves historical execution or establishes physical storage
durability. Executable methods are selected by the caller; packet content
cannot load code.

\paragraph{A module uses a company's existing book.}
Consider the retained return-right example. Its legal book already contains
cash 100, a recovery asset 16 and a refund liability 40. A return module must
refer to those balances. Adding its opening cash or capital to the company
again would create resources that no event supplied.

The finite invocation therefore binds the original entity, legal book,
branch, source context, information view, method and complete opening state.
The trusted return method produces a proposed state transition. It uses the
named book's cash, preserves other legal books and exposes the journal and
gross cash effects. The checked example pays 48 and ends with cash 52. No
commit is implied by the proposal. Changing the owner, method, source, state,
event or account map prevents reuse; missing and false findings remain
distinct from a successful calculation.

The retained owner is \code{Synthetic-return-expected}, with settlement on
30 June 2024. It is not Northstar. Northstar's current language has no
return-right memory, refund/recovery accounts or return-settlement operation.
Its homogeneous inventory and immediate-deductible tax method also cannot
represent this case unchanged. The successful finite proposal identifies a
usable method boundary; native company support needs the state design below.
The separate sector-company method remains report-only and cannot be turned
into a posting operation by relabelling its result.

\paragraph{Returned goods need their own carrying value.}
The original return contract sold ten units at cost 12, with a refund of 20
per returned unit. Its opening estimates recognize a refund liability 40 and
recovery asset 16. These amounts depend on the supplied expected returns,
handling and deterioration. The historical accounting source requires the
refund liability and recovery asset to be distinguished, with the latter
based on former carrying amount less expected recovery costs and loss of
value \citep[IFRS~15.B21--B27]{eu20231803}.

At final settlement, two units return. Their deterioration is 2 per unit,
handling costs 4 per unit and the supplied net realizable value of the lot is
20. Handling is assumed to be selling administration; costs needed to bring
inventory to its location and condition would require a different treatment
\citep[IAS~2.15--16]{eu20231803}. The recovered lot is therefore carried at
$\min\{2(12-2),20\}=20$, and cash payments are 40 for refunds and 8 for handling.

The settlement debits the refund liability 40 and inventory 20, credits the
recovery asset 16 and cash 48, credits cost of sales 4 and debits selling
expense 8. The journal balances and records a book loss of 4. A later sale
of the whole recovered lot for 40 gives revenue 40 and cost of sales 20.
The lot retains its original contract, quantity, carrying amount, tax base
and date so this second event can consume the right stock once.

Simply adding two units to Northstar's homogeneous inventory would value them
at its existing unit cost, which differs from the recovered lot's cost of 10.
The extended book instead reconciles the original homogeneous lot and each
recovered lot separately. Let $q_0$ and $c_0$ denote the homogeneous units and
unit cost, and let $V_\ell$ denote each recovered lot's carrying amount, all
for the same legal book and currency. The proposed inventory control is
\begin{equation}
 I=q_0c_0+\sum_{\ell}V_\ell.
 \label{eq:module-return-inventory}
\end{equation}
The sum preserves different carrying values without silently selecting a
cost-flow method. With no recovered lots, equation~\eqref{eq:module-return-inventory}
reduces to the existing inventory rule. Whole-lot sale is the initial subset;
partial sales and changes of estimates require further specified methods.

\paragraph{Tax history changes the cash needed at the close.}
The same return appears in the joined company case of
Section~\ref{sec:company-composition}. Its legal parent P supplies a complete
trace through settlement, later sale and two tax closes. P's opening, dates
and tax premises are retained; they are not new facts about Northstar.
The case assumes that the original sale and cost were taxed and deducted,
refunds and handling become deductible on settlement, and recovered inventory
enters tax stock at its carrying amount. The supplied tax rate is 25 percent.
The required enactment and recoverability findings remain explicit. Tax bases,
recognition conditions and recovery manner matter independently of book profit
\citep[IAS~12.7--10,24--25,46--51]{eu20231803}.

To see the difference, let $F$ and $R$ be the opening refund liability and
recovery asset, $n$ actual returned units, $u$ refund per unit, $h$ selling
handling per unit and $V$ recovered stock. Under these particular tax premises,
the settlement's book profit $B$ and taxable profit $T$ are
\begin{align}
 B&=F-nu-R+V-nh, & T&=-nu-nh+V,\\
 T-B&=R-F. &&
 \label{eq:module-return-tax}
\end{align}
The tax deduction still includes the refund whose book expense was recognized
before the opening; the book reversal of recovery value has its own earlier
tax difference. Substituting $F=40$, $R=16$, $nu=40$, $nh=8$ and $V=20$ gives
$B=-4$, $T=-28$ and the one-time adjustment $T-B=-24$.

The opening liability's zero tax base gives a deferred tax asset of 10, while
the recovery asset's zero tax base gives a deferred tax liability of 4. Their
reversal adds deferred tax expense 6. P's first-half pretax profit 56 therefore
produces taxable income 32, current tax 8 and after-tax profit 42. Annual
pretax profit 66 gives taxable income 42 and total current tax 10.5; the
second close recognizes only the remaining 2.5. Applying the $-24$ adjustment
afresh to each interval would deduct it twice. The tax history binds the
settlement event and subtracts tax already recognized from the cumulative
calculation.

Multiplying first-half book profit by 25 percent would instead give current
tax 14 and omit the separate deferred expense 6. A coincidentally equal total
expense would not make the current cash payment correct. A complete native
tax close consequently needs tax bases, timing, rate, origin, recovery support
and the full temporary-difference inventory. Without them the tax-dependent
annual report remains unresolved, while earlier valid journals survive as a
dated prefix. The example establishes no actual tax statute.

\paragraph{The statement bridge must retain opening obligations.}
P's annual profit is 49.5, but operating cash is 101.5. Depreciation 20 and
deferred tax expense 6 are noncash charges. Inventory falls by 60, the recovery
asset falls by 16 and the refund liability falls by 40. Another payment is
easy to overlook: the matured opening debt includes 10 of interest expensed
before this year. The selected historical interest policy classifies that
payment as operating, although current finance expense is zero
\citep[IAS~7.31--33]{eu20231803}. The exact indirect bridge is
\begin{equation}
 49.5+20+6+60+16-40-10=101.5.
 \label{eq:module-return-cash}
\end{equation}
Omitting the last term overstates operating cash by 10. Subtracting the whole
debt payment instead would move 90 of financing principal into operating
cash. Equation~\eqref{eq:module-return-cash} agrees with the separate gross
cash events. Opening cash 500, operating cash 101.5, investing cash $-20$ and
financing cash $-90$ yield closing cash 491.5. Opening equity 622 plus profit
49.5 yields 671.5.

These are P's legal-book amounts. Its internal receipt from S remains a legal
cash inflow and is eliminated only in the group report. The same distinction
prevents an apparently balanced module from double-counting internal flows.
Deferred tax is added back once in the indirect bridge; separately adjusting
DTA and DTL again would repeat the same effect.

\paragraph{What the native extension must carry.}
The richer profile needs closed account and component schemas, explicit
return/lot authoring, strict eligibility findings, one chronological shared
journal executor, tax-origin history and extended report mappings. Every
temporary account must enter profit and the closing transfer to retained
earnings. A complete year carries remaining lots, rights, tax bases, event
identities and source context forward. It resets period income and cash-flow
totals without deleting the history needed to prevent a second settlement.
A partial year supplies no complete next-year opening.

The existing company profile remains unchanged. Selecting the richer profile
is explicit; it cannot redefine old record fields or silently replace a
method. The design and exact legal-book trace specify the required repair.
Section~\ref{sec:company-return-reference} now supplies a closed component
authoring profile and a finite reference for return, tax, legal reports and
carry. Section~\ref{sec:company-source-state} adds a finite composed-company
declaration and a separate partial-state description for actual disclosures.
General composition and production integration remain unfinished. Broader
corporate, group, bank and insurance methods retain their separate obligations.

The supporting reference records are under \path{data/accounting-master/}:
\begin{quote}\small
\path{company-portable-authoring-20261023}\\
\path{company-module-boundary-20261023}\\
\path{company-native-return-design-20261023}
\end{quote}
They preserve original bytes, a finite shared-book proposal and the complete
retained legal-book comparison. Independent implementation, domain and human
acceptance remain pending.

\section{From a return declaration to the next annual opening}
\label{sec:company-return-reference}

P's recovered stock gives the language a concrete implementation problem.
The return contract, the legal book and the tax calculation each remember
something different about the same two units. The contract records that the
customer's rights have ended. The stock record keeps the recovered carrying
amount until sale. The tax calculation remembers an adjustment arising from
expense recognized before the year began. A company program must preserve
all three without posting the opening balances again.

The finite reference now connects these records through the complete retained
2025 P trace. It implements return settlement, whole-lot sale, P's selected
current and deferred tax method, legal statements and annual carry. Other
corporate events continue to use the checked retained methods. This division
makes the implementation comparison precise: the new return and tax
calculations can be compared with the earlier company case, while the
surrounding methods supply the same corporate transactions as before.

\paragraph{An author declares facts before choosing an event.}
The versioned \code{ReturnModule}\allowbreak\code{AuthoringDraft/1} describes a component of
one legal book. It identifies the source packet and opening state, entity,
currency, period and account roles. Its contract memory includes sold and
expected-return units, unit refund and cost, estimated recovery costs, and
outstanding refund and recovery balances. The recovered-stock declarations
keep quantities, carrying amounts and tax bases separately. Tax declarations
state the selected timing method, rate and required findings.

These declarations allocate existing accounts. For example, a refund balance
of 40 must agree with the liability in P's opening book; it is not an
instruction to create another liability. The deferred asset and liability
must likewise agree with their declared opening allocations. The remaining
company balances come from the shared opening, rather than from a second
component wallet.

Typed records and text declarations contain the same fields. The text form
uses named declarations for the component, lots, tax premises and dated
settlement or sale. It contains no executable Python, external include or
opaque JSON body. Both forms preserve canonical exact rational amounts, strict
Boolean findings and source references. Rendering a typed declaration as text
and reading it again must reproduce the same value. This is a matched
authoring comparison with the typed-library baseline, not evidence that text
is preferable for every user.

Unknown values are explicit. Suppose the tax rate has not been established
but the author also supplies a refund liability of 41 while two expected
returns at 20 require 40. The contradiction is already decidable and the
record is invalid. The unknown tax rate cannot conceal it. If the amounts
agree but a required recovery finding is false, the selected method is
unsupported. If the required finding is unknown, the affected calculation
remains unresolved. Arithmetic checks inspect their own dependencies rather
than replacing unknowns with zero or postponing every known check.

Source correspondence is another question. A syntactically valid source
pointer may name the wrong disclosure. The retained example therefore checks
the entire declared projection against P's original packet, including owner,
dates, amounts and opening identity. Changing P to Northstar fails that
correspondence even though both names are valid text. A correctly bound local
packet still needs separate evidence for external authenticity and any
historical admission claim.

\paragraph{One event changes one shared book.}
The reference state combines P's account balances with return memory,
recovered lots, dated tax adjustments, classified gross cash and lifetime
event identities. Every account belongs to one financial-statement class;
every temporary account enters profit. A proposed return first checks its
component, quantities, findings and date. It then derives the settlement
journal already shown in Section~\ref{sec:company-module-language}, creates
the separate recovered lot and closes the return rights. The surrounding
inventory allocation stays intact.

The proposal identifies the complete preceding state, the exact event and
the selected method bundle. Applying it repeats the calculation with the
trusted local method and compares the journal, gross cash and complete
after-state. A changed cash amount cannot become valid merely by supplying a
new hash. A stale proposal, repeated settlement or sale of an already disposed
lot is refused. This is an in-process reference for checking a transition;
durable publication and concurrent writers require the separate storage design.

Cash sufficiency is checked after deriving the complete proposed journal.
If P cannot pay the refund and handling costs, the event returns a funding
failure and leaves the preceding book intact. It does not create financing.
Missing tax evidence has a different consequence: the valid non-tax events
can remain as an exact dated prefix, but the tax-dependent annual statement
and next-year opening remain unavailable. Thus a useful partial history
survives without being presented as a complete year.

The retained adapter has an explicit, narrower role. It recomputes the other
corporate events from their original state and checked methods, verifies that
their P opening balances agree with the shared book, and supplies their
journals and cash origins. Saved expected journals are comparison evidence,
not instructions to the executor. S's book and the group ownership state
remain in that retained method. Replaying them checks compatibility; it does
not claim a second independent implementation of those same corporate rules.

\paragraph{Tax accumulates within a year and resets at its boundary.}
The settlement records its taxable-minus-book adjustment of $-24$ with its
event identity and date. At either tax close, let $B_y(t)$ be P's cumulative
pretax book income from the start of year $y$ through cutoff $t$, let
$A_y(t)$ be the sum of the selected method's dated tax adjustments in that
interval, and let $C_y(t^-)$ be current tax already recognized in the year.
Under the retained constant rate $r=1/4$, the additional current tax is
\begin{equation}
 \Delta C_y(t)=r\{B_y(t)+A_y(t)\}-C_y(t^-).
 \label{eq:return-reference-current-tax}
\end{equation}
The final subtraction prevents an interim tax charge from being paid again.
At the first close, $B_y=56$, $A_y=-24$ and $C_y(t^-)=0$, giving 8. At the
annual close, $B_y=66$ and the same cumulative adjustment remains $-24$;
subtracting the existing 8 gives the remaining 2.5. Negative additional tax
would require a refund or other tax method that this finite reference does
not supply. The opening deferred balances reverse separately, producing
deferred expense 6.

The annual legal statements reconcile profit 49.5, cash 491.5 and equity
671.5. Their operating cash of 101.5 agrees with both the gross payments and
the indirect bridge in equation~\eqref{eq:module-return-cash}. In particular,
the opening interest payment remains a cash adjustment even when current
finance expense is zero. Legal receipts from S stay in P's cash flow; they
are eliminated only by the separate group method.

At year end, the reference transfers profit to retained earnings and clears
the income and period cash accumulators. It retains the source context,
settled return component, disposed-lot identity and lifetime event history.
The old tax adjustment remains available to explain 2025, but its date is
outside $A_{2026}(t)$. In a deliberately inactive following year, equation~
\eqref{eq:return-reference-current-tax} therefore gives zero additional tax
and the deferred reversal is not repeated. This checks annual memory; it is
not a prediction of an inactive business.

The complete reference compares every P account and gross cash origin after
all thirteen dated transitions, then checks the annual statements and carry.
It also preserves the existing six company-authoring paths. The next
section exposes the selected methods, source state and schedule in a finite
composed-company declaration, then distinguishes the partial inputs available
from actual disclosures. General changes
in return estimates, partial lot sales, other tax treatments and production
execution still need their own design and conformance cases.

The matched authoring and execution records are under
\path{data/accounting-master/}:
\begin{quote}\small
\path{company-return-schema-20261023}\\
\path{company-return-execution-20261024}
\end{quote}
They retain
the original source comparisons and failure cases. The resulting methods are
conditional finite references; independent implementation, domain review and
human acceptance of the governing design remain pending.

\section{A company program begins with the state it can justify}
\label{sec:company-source-state}

P's return example has an opening book, identified lots, dated events and
explicit tax premises. That information allows the reference calculation to
reach an annual result. An actual company's published statements usually
offer a different starting point: amounts reported for a particular group
and period, together with policies and notes that leave some underlying
contracts and allocations unknown. The language must accommodate both
starting points. An author should be able to select methods for a complete
conditional example and also describe a source-constrained partial state
whose missing facts still prevent those methods from running.

\subsection{Selecting methods for a complete retained company}

For P, the finite composed-company declaration names the source packet and
opening, legal-book context, reporting period, return component and all
thirteen events in their original order. Each event selects a declared
capability. Return settlement and recovered-lot sale use the native return
method; the two tax closes use the native tax method; the other transactions
use the previously checked corporate methods. Annual reporting and carry
have their own explicit selections. The caller provides the available
implementations and their identities. A capability name written in a
program cannot import arbitrary code, and an unavailable capability cannot
silently select a different method.

Matched text and typed records express the same declaration. The relevant
draft is \path{RetainedCompanyProgramDraft/1}; its example and evaluator
are identified at the end of this section. Binding the program to the
retained case means comparing its full source, opening, event schedule and
method selections, rather than accepting a familiar company name as proof
of identity. An author cannot rename P as Apple and retain the old source
facts. Nor does a correctly hashed result establish its own correctness:
the reference verifier recomputes the proposals, annual report and carry.

Review findings remain separate from immutable source facts. If a required
tax finding is unknown, the first seven valid events remain available,
the six subsequent events remain unexecuted, and no complete annual report
or next opening is returned. A negative finding instead records that the
selected treatment is unsupported. Editing a finding changes the authored
program's identity; it does not edit the source packet. With the retained
findings supplied, the composed calculation reproduces P's profit of 49.5,
cash of 491.5, equity of 671.5 and operating cash flow of 101.5. This
completes the finite declaration for the case while leaving broader company
composition and method admission to their own specifications.

\subsection{What the disclosed balance sheet leaves unresolved}

Consider the retained Apple balance sheet at 28 September 2024. It reports
cash and cash equivalents of 29,943, inventory of 7,286 and net property,
plant and equipment of 45,680, in millions of US dollars
\citep[p.~31]{apple2024}. These are useful observations. They do not, however,
identify the unrestricted cash available to each legal entity, inventory
quantities and lot costs, or the original costs, accumulated depreciation,
service dates and remaining lives of asset cohorts. Supplying one invented
unit cost or useful life would choose an economic history that the
disclosure has not established.

The input-readiness reference therefore preserves the original statements
and asks what each selected execution method still requires. It keeps all
139 package-inventory entries, the 102 annual and 100 interim readings,
and the two collections of 27 opening or comparative balances. Fifteen
dependencies connect those records and the original note obligations to
company inputs, affected operations, reports and annual carry. Missing
claims, asset cohorts, tax bases and dated events remain visible beside the
known displayed amounts.

There are also restrictions in the selected methods themselves. The
operations draft requires a full calendar year, a fixed account registry,
homogeneous inventory cost, specified asset cohorts and a simple tax
treatment. Apple's retained consolidated disclosures do not satisfy those
requirements merely because their balance sheet balances. P's composed
program has a different restriction: it binds exactly the retained
synthetic case. The readiness result keeps missing company state and
unsupported method scope as separate reasons for refusing execution.
That refusal concerns the selected profiles; it does not establish that
Apple cannot be represented by a broader model.

Known arithmetic can still be retained without pretending to solve those
problems. The displayed current and noncurrent term-debt captions sum to
\(10{,}912+85{,}750=96{,}662\) at the annual date. This is a sum of two
reported captions, not total financing obligations or an allocation to
individual instruments. The source's dash for annual debt-issuance cash
flow remains a dash. Unknown original-byte publication and ingestion clocks
limit historical as-of admission, but they do not prevent reading the
named retained source. The unresolved actual-company history and the earlier
conditional subsystem calculations remain separate, as before.

\subsection{Fiscal dates belong to the source}

A partial-state declaration can carry the reported measurements without
manufacturing an executable opening. The draft
\path{SourceConstrainedCompanyStateDraft/1} binds the exact readiness request
and result, declares its fiscal intervals and names each selected observation.
It includes all 54 current annual and interim balance captions. Each binding
keeps the original measurement date, consolidated perimeter, basis, units,
precision and source-admission status. Prior comparisons and the remaining
readings stay available in the referenced readiness result; they have not
all become executable state constraints.

The retained annual source covers 1 October 2023 through 28 September 2024;
the following interim source covers 29 September through 28 December 2024
\citep{apple2024,apple2025q1}. Using the half-open interval convention already
introduced in Chapter~\ref{ch:history}, these become
\begin{equation}
\label{eq:source-fiscal-periods}
  [2023\text{-}10\text{-}01,\,2024\text{-}09\text{-}29),
  \qquad
  [2024\text{-}09\text{-}29,\,2024\text{-}12\text{-}29).
\end{equation}
The exclusive end is the day after the source's inclusive closing date.
The first interval has 364 days and the second 91; they meet without a
gap or overlap. A balance-sheet observation dated 28 September remains at
that date, rather than moving to the exclusive boundary on 29 September.

Those date calculations specify reporting coverage. They do not select an
interest day-count convention, a depreciation fraction or tax-year law.
The generic structural validator admits explicit non-calendar periods;
the retained-source binder separately rejects a period that disagrees with
the original observation. Earlier fiscal-closing examples already permit
short periods. The addition here makes fiscal intervals directly authorable
alongside a partial source state, while preserving the stricter scope of
the existing execution methods.

\subsection{A rounding assumption constrains a measurement}

The default interpretation of a displayed amount is deliberately modest:
the source displays this value, while its unrounded underlying measurement
is unresolved. Storing the displayed integer as an exact rational preserves
the token; it does not establish exact equality with the company's
underlying balance. Even an interval requires a rounding convention that
the retained precision annotation does not by itself establish.

An author may supply a separate rounding hypothesis. Suppose the reporting
quantum is \(q>0\), the displayed value \(d\) is an integer multiple of
\(q\), and the rule rounds to the nearest multiple with ties away from
zero. Let \(x\) denote the underlying value of the same reported
measurement: it has the same date, perimeter, unit and basis, not a newly
inferred legal-account meaning. The midpoints adjacent to \(d\) are
\(d-q/2\) and \(d+q/2\). At a midpoint the value with the larger absolute
magnitude wins. It follows that the values producing \(d\) are
\begin{equation}
\label{eq:source-rounding-preimage}
  x\in
  \begin{cases}
    [d-q/2,\,d+q/2), & d>0,\\
    (d-q/2,\,d+q/2], & d<0,\\
    (-q/2,\,q/2), & d=0.
  \end{cases}
\end{equation}
For positive \(d\), the lower tie rounds up to \(d\), while the upper tie
rounds away to \(d+q\). For negative \(d\), the endpoint inclusion
reverses. Neither midpoint rounds to zero, which explains the two open
endpoints in the last case.

Under that explicit hypothesis with \(q=1\) million dollars, the annual
cash-and-equivalents caption of 29,943 permits
\([29{,}942.5,29{,}943.5)\). The displayed retained deficit of
\(-19{,}154\) permits \((-19{,}154.5,-19{,}153.5]\)
\citep[p.~31]{apple2024}. These are conditional deductions, not evidence
that Apple uses the proposed tie-breaking rule. The assumption has its own
identity and provenance, and it remains distinct from accepted source facts.
Without it, the declaration retains the displayed value and leaves the
underlying measurement unresolved.

The intervals neither allocate cash to legal entities nor split inventory
into lots. They are not assumed independent, and their Cartesian product
is not asserted to be a feasible company history. Every missing-state
dependency therefore survives the rounding calculation. The reference
returns a partial description, with joint feasibility and actual history
unknown, and supplies no executable opening. This lets the language retain
what a source says while making the next evidence or modelling choice
precise enough to review.

The closed typed records, named text examples and bounded reference checks
are under \path{data/accounting-master/}:
\begin{quote}\small
\path{company-composed-profile-20261024}\\
\path{real-company-input-readiness-20261024}\\
\path{source-constrained-state-20261024}
\end{quote}
Their checks cover source identity, complete dependency retention, matching
text and typed interpretations, fiscal boundaries and signed rounding ties.
They reject changed source dates, legal-book relabelling, invented exact
balances and a forged executable result. The next section extends the source
description beyond current balance captions while preserving the difference
between displayed relations and latent company state. General source binders, admissible
state reconstruction, broader accounting methods and independent acceptance
of the governing design remain unfinished.

\input{sections/statement-context-relations}

\input{sections/joint-source-measurements}

\input{sections/temporal-statement-network}

\input{sections/conditional-history-families}

\input{sections/typed-economic-history}

\input{sections/conditional-method-families}



